
Stand der Technik

Am Anfang steht die Frage nach den Gemeinsamkeiten aller verwendeten Daten.
Dazu ein Beispiel in der Syntax der Programmiersprache Java:

int i = 0;

Um einen Wert, der im Arbeitsspeicher liegt, nutzen zu können,
muss man ihn referenzieren können.
Dies geschieht durch die Variable mit dem "Namen" i.

Außerdem wird (bei typisierten Sprachen) die Größe
des im Arbeitsspeicher zu reservierenden Bereiches
durch einen "Typ" vorgegeben.

Schließlich folgt der eigentliche, initial als Literal
zugewiesene "Wert" 0.

TODO: Recherche dynamische Typisierung z. B. in JavaScript?
a = 3;
Wie intern gespeichert? Als String?

Man hat also ein Schlüssel-Wert-Paar:
- der Variablenname i steht für den Schlüssel;
- die Zahl 0 im Beispiel für den Wert.

In typisierten Sprachen wird dieses Paar ergänzt durch einen Typ.

Damit lässt sich vorerst zusammenfassen, dass Daten beschrieben
werden durch:
- Name
- Typ
- Wert

Neben primitiven Werten können auch komplexe verwendet werden, z. B.:

Address a = new Address();

Da es sich bei "Address" um eine zusammengesetzte Datenstruktur
und keinen primitiven Wert handelt, bietet es sich an, den
neutraleren Begriff "Modell" statt "Wert" zu verwenden:

TABELLE:

Bezeichnung     Beschreibung/Synonyme
Name (name)     Schlüssel (key), Bezeichner (variable)
Typ (type)      Typ, Art der Abstraktion
Modell (model)  Wert (value)
